\section{Survey Methodology}
\label{sec:method}

\subsection{Research questions}
\begin{enumerate}[label=\textbf{RQ\arabic*.},leftmargin=2.6em]
  \item Which engineering-information capabilities do industrial AI agents require to
  ground observations and actions?
  \item Which standards, information models, semantic technologies and knowledge
  representations provide them?
  \item Where do current industrial-agent architectures obtain identity, units, topology,
  behaviour, constraints, provenance and runtime state?
  \item Which capabilities remain missing or weakly supported?
  \item Which architectural patterns combine LLM reasoning with deterministic engineering
  validation?
  \item What evidence shows that engineering semantics improve grounding, action
  correctness, transferability, auditability or context efficiency?
\end{enumerate}
RQ1 is answered by the AUF derivation (Section~\ref{sec:auf}), RQ2 by the technology
landscape and matrix (Sections~\ref{sec:landscape}--\ref{sec:capabilities}), RQ3--RQ5 by
the agent-literature coding (Sections~\ref{sec:agents}--\ref{sec:architecture}), and RQ6
by the reported evaluations and the case study (Section~\ref{sec:case}).

\subsection{Sources and search strategy}
Four evidence streams are distinguished and never pooled. \emph{Specifications} (IEC,
ISO, ISA, VDI/VDE/NAMUR, W3C, ETSI, OPC Foundation, IDTA, ECLASS, GS1, IDSA) are collected
from the issuing bodies and scored on what they standardise. \emph{Peer-reviewed
literature} published 2021--2026 is searched in IEEE Xplore, Scopus, Web of Science, ACM
Digital Library, ScienceDirect and SpringerLink by the authors, and in OpenAlex and
Crossref by released tooling; Crossref is relevance-ranked and capped, so it supplements
rather than replaces the boolean database searches. \emph{Preprints} (arXiv, posted
content) are coded but labelled non-peer-reviewed. \emph{Vendor material} supports only
implementation-capability statements. Foundational earlier work (retrieval augmentation,
tool use, constraint languages) enters through backward snowballing. Rather than one
query, the protocol defines eight topical families, (F1) LLM agents in industrial
automation and cyber-physical systems, (F2) industrial knowledge graphs and ontologies with
LLMs, (F3) \opcua{}, AAS and AutomationML with LLMs or agents, (F4) runtime validation and
guardrails in control systems, (F5) semantic grounding and interoperability, (F6) Graph- and
KG-RAG on engineering data, (F7) neurosymbolic, ontology-grounded and typed-tool agents, and
(F8) safety, provenance and trust, plus a family (R1) dedicated to related surveys. Exact
database-specific strings are frozen in protocol v\SurveyProtocolVersion{}, released with
the article; forward and backward snowballing from included studies complements them.

\subsection{Eligibility, screening and quality}
Studies are included if they address industrial AI or LLM agents or the engineering
semantics they consume (I1), industrial information models or knowledge graphs with an AI
consumer (I2), grounding, tool use, retrieval or runtime action validation in an
engineering context (I3), or constraints, provenance or state models used to check AI
outputs (I4). Exclusion codes are: no industrial context (E1); generic RAG or graph
construction without engineering semantics (E2); no agent or LLM (E3); duplicate or
superseded version (E4); insufficient documentation (E5); inaccessible full text (E6);
language or period (E7). Surveys (I5) feed only the related-survey gap analysis. Records are
deduplicated by DOI and by normalised title, year
and first author; near-identical titles are listed for manual resolution rather than
merged. A deterministic, versioned automation rule then marks records ineligible unless
title or abstract contains both an AI/semantics term and an industrial term; every
excluded record and the failing rule are released, as PRISMA~2020 requires for automation
tools. Title/abstract and full-text screening use spreadsheets that preserve each
decision and reason. A seeded random 20\% sample is designated for double coding at both
stages. Five quality items (architecture described, evaluation method, artefact
availability, isolation of the semantic contribution, safety limitations) are scored 0--2
and reported, not used for exclusion.

\subsection{Coding and agreement}
Each included system is coded for grounding category (Section~\ref{sec:agents}), retrieval
mechanism, tool interface, representation, each AUF capability on a 0--3 scale (absent,
implicit, explicit, checked at runtime), runtime validation, controller authority and
evaluation level (none, simulation, testbed, physical). Technologies are coded with the
rubric in Table~\ref{tab:rubric}, and each score is linked to a specification reference.
Agreement on the double-coded sample is reported as quadratic-weighted Cohen's $\kappa$
and ordinal Krippendorff's $\alpha$ with adjudicated disagreements; at the time of writing
the second rater's coding is pending ($\kappa$: \AufKappa{}, $\alpha$: \AufAlpha{}), and all
scores are first-rater codes.

\subsection{Capability levels}
Every capability statement records its level: \emph{specification} capability (what a
standard can express by default), \emph{implementation} capability (what a cited
implementation or profile exposes), and \emph{runtime} availability (what the agent
actually receives or checks when deciding). Matrix A (Section~\ref{sec:capabilities}) codes
specification capability; Matrix B (Section~\ref{sec:agents}) codes runtime availability
and records ``not reported'' rather than inferring capabilities a paper does not
demonstrate.

\subsection{Search status}
Figure~\ref{fig:prisma} is generated from the screening sheets. The open-source harvest
(\SurveyHarvestDate{}) identified \SurveyIdentified{} records (\SurveyCrossref{} Crossref,
\SurveyArxiv{} arXiv, \SurveyOpenAlex{} OpenAlex; incomplete for: \SurveyFailedSources{}).
After removing \SurveyDuplicates{} duplicates and marking \SurveyAutoExcluded{} records
ineligible by the automation rule, \SurveyUnique{} records await title/abstract
screening, and \SurveyRelatedCandidates{} industrial survey candidates await full-text
verification for Table~\ref{tab:related-surveys}. Subscription-database exports are
imported with the same tooling, after which all counts are regenerated.

\begin{figure}[t]
\centering
\resizebox{\linewidth}{!}{\begin{tikzpicture}[node distance=3.2mm and 3mm, >=Stealth]
\node[draw, rounded corners=1pt, align=center, font=\scriptsize, text width=3.6cm, minimum height=7mm] (id) {Records identified: 17197\\{\tiny arxiv 1186, crossref 9000, openalex 7011}};
\node[draw, align=center, font=\scriptsize, text width=2.9cm, minimum height=7mm, right=of id] (dup) {Duplicates removed: 3797\\Ineligible by automation: 6373};
\node[draw, rounded corners=1pt, align=center, font=\scriptsize, text width=3.6cm, minimum height=7mm, below=of id] (scr) {Records screened: 7027\\\textit{screening in progress}};
\node[draw, align=center, font=\scriptsize, text width=2.9cm, minimum height=7mm, right=of scr] (tax) {Excluded (title/abstract): 0\\pending: 7027};
\node[draw, rounded corners=1pt, align=center, font=\scriptsize, text width=3.6cm, minimum height=7mm, below=of scr] (ft) {Full texts sought: 0\\not retrieved: 0};
\node[draw, align=center, font=\scriptsize, text width=2.9cm, minimum height=7mm, right=of ft] (ftx) {Excluded (full text): 0\\{\tiny --}};
\node[draw, rounded corners=1pt, align=center, font=\scriptsize, text width=3.6cm, minimum height=7mm, below=of ft] (inc) {Included and AUF-coded: 0\\pending: 0};
\draw[->] (id) -- (scr); \draw[->] (scr) -- (ft); \draw[->] (ft) -- (inc);
\draw[->] (id) -- (dup); \draw[->] (scr) -- (tax); \draw[->] (ft) -- (ftx);
\end{tikzpicture}
}
\caption{PRISMA 2020 flow generated from the screening sheets; counts update automatically
as screening proceeds.}
\label{fig:prisma}
\end{figure}
